\documentclass[11pt,a4paper]{article}

% ---------------------------------------------------------------- idioma ----
\usepackage[utf8]{inputenc}
\usepackage[T1]{fontenc}
\IfFileExists{lmodern.sty}{\usepackage{lmodern}}{}
\IfFileExists{spanish.ldf}{%
  \usepackage[spanish,es-nodecimaldot,es-tabla]{babel}%
}{%
  % (solo por si el compilador no tiene babel-spanish; en Overleaf/TeX Live completo se usa la rama de arriba)
  \usepackage[english]{babel}%
  \renewcommand{\contentsname}{Índice}\renewcommand{\figurename}{Figura}%
  \renewcommand{\tablename}{Tabla}\renewcommand{\refname}{Referencias}%
}

% ------------------------------------------------------------- paquetes -----
\usepackage{amsmath,amssymb,bm}
\usepackage{graphicx}
\usepackage{booktabs}
\usepackage[margin=2.6cm]{geometry}
\usepackage{xcolor}
\usepackage{caption}
\usepackage{enumitem}
\usepackage{listings}
\usepackage[colorlinks=true,linkcolor=blue!50!black,urlcolor=blue!60!black,citecolor=blue!50!black]{hyperref}
\captionsetup{font=small,labelfont=bf}
\setlength{\parskip}{3pt}

% =================== DATOS QUE DEBES PERSONALIZAR ===========================
\newcommand{\autor}{Nombre Apellidos}
\newcommand{\repoURL}{https://github.com/TU\_USUARIO/valle-parabolico}      % <- cambia esto
\newcommand{\repoHREF}{https://github.com/TU_USUARIO/valle-parabolico}     % <- y esto (misma URL)
% ============================================================================

% --------------------------------------------------------------- macros ------
\newcommand{\dd}{\mathrm{d}}
\newcommand{\pd}[2]{\dfrac{\partial #1}{\partial #2}}
\newcommand{\ddt}[1]{\dfrac{\dd #1}{\dd t}}
\newcommand{\ve}[1]{\vec{#1}}
\DeclareMathOperator{\sgn}{sgn}

\lstset{
  basicstyle=\ttfamily\footnotesize, numbers=left, numberstyle=\tiny\color{gray},
  breaklines=true, frame=single, columns=fullflexible, keepspaces=true,
  keywordstyle=\color{blue!60!black}, commentstyle=\color{green!40!black}\itshape,
  stringstyle=\color{red!50!black}, language=Python, showstringspaces=false,
  inputencoding=utf8, extendedchars=true,
  literate={á}{{\'a}}1 {é}{{\'e}}1 {í}{{\'i}}1 {ó}{{\'o}}1 {ú}{{\'u}}1 {ñ}{{\~n}}1
           {Á}{{\'A}}1 {É}{{\'E}}1 {Í}{{\'I}}1 {Ó}{{\'O}}1 {Ú}{{\'U}}1
           {²}{{$^2$}}1 {·}{{$\cdot$}}1 {π}{{$\pi$}}1 {ª}{{$^{\underline{a}}$}}1
}

\title{\vspace{-1cm}\textbf{El Valle Parabólico}\\[2mm]
\large Dinámica lagrangiana de un vehículo sobre rails parabólicos fijos y giratorios\\[1mm]
\normalsize Ejercicio 1 de la convocatoria extraordinaria de julio/septiembre (curso 2025--26)}
\author{\autor\\[1mm] \small Mecánica y Ondas I --- Grado en Física, Universidad de Córdoba}
\date{\today}

\begin{document}
\maketitle

\begin{abstract}
Se estudia el movimiento de un vehículo (masa puntual, sin rozamiento) sobre unos rails con forma
de parábola vertical $z=ax^2$ mediante la formulación lagrangiana. Se resuelven los tres casos del enunciado:
(b) rails fijos, un grado de libertad; (c) rails que giran alrededor del eje vertical, dos grados de libertad y
conservación del momento angular; y (d) cálculo de la reacción normal mediante una lagrangiana auxiliar con
multiplicador de Lagrange. Además de las soluciones analíticas (periodo exacto mediante una integral elíptica,
potencial efectivo, órbitas circulares, ángulo apsidal, reacción normal máxima $N_{\max}=mg(1+4a^2x_0^2)$)
se ha programado una simulación numérica que las reproduce con errores relativos
inferiores a $10^{-7}$, y se hace un estudio sistemático de las condiciones iniciales, de los casos límite y de la influencia de los
parámetros del sistema. El código y las animaciones están disponibles en
\href{\repoHREF}{\texttt{\repoURL}}.
\end{abstract}

\tableofcontents

% =============================================================================
\section{Introducción y enunciado}
% =============================================================================
El ejercicio propone tres atracciones de feria. Este trabajo resuelve la \emph{primera} (ejercicio~1,
«Formulación lagrangiana»), en sus variantes de rails fijos y giratorios. Responde a las siguientes cuestiones:
\begin{enumerate}[label=\alph*),leftmargin=*]
  \item Deducir las ecuaciones de Lagrange para campos conservativos a partir de las de Newton (principio de
        d'Alembert, desplazamiento virtual, grados de libertad, ligaduras, coordenadas generalizadas y relación
        lógica entre ambas formulaciones). \emph{Resumen en la sección~\ref{sec:teoria}; deducción completa en el apéndice~\ref{app:lagrange}.}
  \item Vehículo sobre rails parabólicos fijos: nombre de la atracción, grados de libertad, coordenada
        generalizada, lagrangiana, constantes del movimiento y ecuaciones del movimiento (sección~\ref{sec:b}).
  \item Rails que giran alrededor del eje vertical: grados de libertad, coordenadas, lagrangiana, constantes del
        movimiento y ecuaciones del movimiento (sección~\ref{sec:c}).
  \item Rails fijos: lagrangiana auxiliar para obtener la reacción normal, ecuaciones, número de incógnitas y
        significado de la magnitud que representa la reacción (sección~\ref{sec:d}).
\end{enumerate}

\paragraph{Modelo y convenios.} El vehículo y sus ocupantes se aproximan por una masa puntual $m$; se desprecian el rozamiento con los rails y con el aire.
Los rails son la curva $z=ax^2$ ($a>0$, unidades m$^{-1}$) contenida en un plano vertical; el eje $z$ apunta hacia arriba y
la gravedad es $\vec g=-g\,\hat z$. Para los resultados numéricos se toma
\begin{equation}
 a=0.1\ \mathrm{m^{-1}},\qquad g=9.81\ \mathrm{m/s^2},
\end{equation}
que corresponde a una atracción realista (a 10~m del eje los rails ya han subido 10~m, y el radio de curvatura en el fondo es $R=1/(2a)=5$~m).
La masa se cancela en todas las ecuaciones del movimiento (principio de equivalencia), por lo que las energías se dan por unidad de masa cuando se comparan resultados.
Todas las figuras y números de este documento se generan con el código del repositorio (sección~\ref{sec:codigo}).

% =============================================================================
\section{Resumen teórico: de Newton a Lagrange (apartado a)}\label{sec:teoria}
% =============================================================================
El desarrollo completo está en el apéndice~\ref{app:lagrange}; aquí se recoge la cadena lógica.

\begin{enumerate}[label=\roman*),leftmargin=*]
\item \textbf{Principio de d'Alembert.} Para $N$ partículas, $m_i\ddot{\vec r}_i=\vec F_i^{(a)}+\vec f_i$ (fuerzas aplicadas y de ligadura).
Un \emph{desplazamiento virtual} $\delta\vec r_i$ es un cambio infinitesimal de la configuración, \emph{instantáneo} ($\delta t=0$) y compatible con
las ligaduras en ese instante. Para ligaduras ideales (como los rails lisos) las fuerzas de ligadura no realizan trabajo virtual,
$\sum_i\vec f_i\cdot\delta\vec r_i=0$, y queda
\begin{equation}
 \sum_i\left(\vec F_i^{(a)}-m_i\ddot{\vec r}_i\right)\cdot\delta\vec r_i=0 .
\end{equation}
\item \textbf{Grados de libertad y coordenadas generalizadas.} Con $k$ ligaduras holónomas $f_\alpha(\vec r_1,\dots,\vec r_N,t)=0$ el sistema
tiene $n=3N-k$ grados de libertad, y se pueden elegir $n$ coordenadas independientes $q_j$ tales que $\vec r_i=\vec r_i(q_1,\dots,q_n,t)$ cumple automáticamente las ligaduras.
\item \textbf{Ecuaciones de Lagrange.} Al escribir $\delta\vec r_i=\sum_j\frac{\partial\vec r_i}{\partial q_j}\delta q_j$, y ser los $\delta q_j$ independientes:
\begin{equation}
 \ddt{}\!\left(\pd{T}{\dot q_j}\right)-\pd{T}{q_j}=Q_j,\qquad Q_j=\sum_i\vec F_i^{(a)}\cdot\pd{\vec r_i}{q_j}.
\end{equation}
Si las fuerzas son conservativas, $\vec F_i=-\nabla_iV$, entonces $Q_j=-\partial V/\partial q_j$ y, con $L=T-V$,
\begin{equation}
 \boxed{\ \ddt{}\!\left(\pd{L}{\dot q_j}\right)-\pd{L}{q_j}=0,\qquad j=1,\dots,n .\ }
\end{equation}
\item \textbf{Dependencia lógica.} Newton $+$ ligaduras ideales $\Rightarrow$ d'Alembert $\Rightarrow$ Lagrange. Cada paso es reversible
(los $\delta q_j$ son arbitrarios y la matriz $\partial\vec r_i/\partial q_j$ tiene rango máximo), de modo que las ecuaciones de Lagrange equivalen a las de Newton
\emph{proyectadas sobre las direcciones permitidas por las ligaduras}. La formulación lagrangiana no contiene nada nuevo desde el punto de vista del contenido físico:
su ventaja es que elimina las fuerzas de ligadura (que dejan de ser incógnitas) y permite usar coordenadas adaptadas al problema y detectar simetrías.
\emph{Precio a pagar:} las fuerzas de ligadura no se obtienen, salvo que se recuperen con multiplicadores de Lagrange (apartado d).
\end{enumerate}

% =============================================================================
\section{Apartado b): atracción fija «El Valle Parabólico»}\label{sec:b}
% =============================================================================
\textbf{Nombre propuesto:} \emph{El Valle Parabólico}.

\paragraph{Grados de libertad y coordenada generalizada.}
La partícula se mueve en el plano vertical $xz$ (dos coordenadas, $y=0$) con una ligadura holónoma, esclerónoma y bilateral,
$f(x,z)=z-ax^2=0$. Por tanto hay $2-1=\mathbf{1}$ \textbf{grado de libertad}. Se elige como coordenada generalizada la abscisa
$q=x$; la posición es $\vec r=(x,\,0,\,ax^2)$ y $\dot z=2ax\dot x$.

\paragraph{Lagrangiana.}
Con $T=\tfrac12m(\dot x^2+\dot z^2)$ y $V=mgz$:
\begin{equation}\label{eq:Lb}
 L_b(x,\dot x)=\tfrac12m\left(1+4a^2x^2\right)\dot x^2-mgax^2 .
\end{equation}

\paragraph{Constantes del movimiento.}
$L_b$ no depende explícitamente del tiempo, luego se conserva la función de Jacobi $h=\dot x\,\partial L/\partial\dot x-L$. Como la
energía cinética es cuadrática en $\dot x$ y la ligadura no depende de $t$, $h$ coincide con la \textbf{energía mecánica}:
\begin{equation}
 E=\tfrac12m\left(1+4a^2x^2\right)\dot x^2+mgax^2=\text{cte}.
\end{equation}
La coordenada $x$ no es cíclica ($\partial L/\partial x\neq0$), así que no hay más integrales primeras.
(Con un solo grado de libertad la conservación de $E$ basta para integrar el problema por cuadraturas.)

\paragraph{Ecuación del movimiento.}
\begin{equation}\label{eq:movb}
 \boxed{\ \left(1+4a^2x^2\right)\ddot x+4a^2x\,\dot x^2+2ga\,x=0\ }
\end{equation}
(deducción en el apéndice~\ref{app:b}). Obsérvese que $m$ no aparece.

\paragraph{Pequeñas oscilaciones y periodo exacto.}
Para $|2ax|\ll1$ la ecuación se reduce a $\ddot x+2ga\,x=0$: oscilaciones armónicas con
\begin{equation}
 \omega_0=\sqrt{2ga},\qquad T_0=\frac{2\pi}{\sqrt{2ga}}=4.486\ \mathrm{s}\quad(a=0.1\ \mathrm{m^{-1}}).
\end{equation}
Para amplitud arbitraria $A$ (punto de retorno, $gaA^2=E/m$) el periodo se obtiene exactamente de la conservación de la energía:
\begin{equation}\label{eq:Tb}
 T(A)=\frac{4}{\sqrt{2ga}}\;\mathcal{E}\!\left(-4a^2A^2\right)\simeq T_0\left(1+a^2A^2-\tfrac34a^4A^4+\dots\right),
\end{equation}
donde $\mathcal E(m)=\int_0^{\pi/2}\sqrt{1-m\sin^2\theta}\,\dd\theta$ es la integral elíptica completa de segunda especie
(apéndice~\ref{app:b}). Para amplitudes grandes ($aA\gg1$), $T\to4\sqrt{2a/g}\,A$: crece linealmente con $A$.
A diferencia de un oscilador armónico, \emph{el movimiento no es isócrono}: cuanto mayor es la amplitud, más lento es.

\subsection{Estudio numérico}
La ecuación \eqref{eq:movb} se integra con un método Runge--Kutta de orden 8 (DOP853, tolerancias $10^{-11}$).
La energía se conserva con un error relativo máximo de $3\times10^{-10}$ en 200~s (figura~\ref{fig:benergia}, izquierda).

\begin{figure}[ht]\centering
\includegraphics[width=\textwidth]{figuras/b_trayectorias.pdf}
\caption{Rails fijos, $a=0.1$~m$^{-1}$. Izquierda: $x(t)$ soltando desde el reposo en distintas posiciones $x_0$.
Derecha: retratos de fases, curvas cerradas de energía constante: el movimiento es siempre periódico, sin separatrices.}
\label{fig:btray}
\end{figure}

\paragraph{Condiciones iniciales.} Soltando desde el reposo en $x_0$ la amplitud es $A=x_0$ (figura~\ref{fig:btray}). Si se lanza con velocidad $v_0$ desde el fondo
($x_0=0$), la energía da $x_{\max}=v_0/\sqrt{2ga}$, que coincide con la simulación (tabla~\ref{tab:b}, derecha); el caso $x_0=0,\ v_0=0$ es el \textbf{equilibrio estable}.
No existe ninguna condición inicial con movimiento no acotado: el pozo de potencial $mgax^2$ crece sin límite.

\begin{table}[ht]\centering
\resizebox{\textwidth}{!}{%
\begin{tabular}{rrrrr@{\qquad}rrr}
\toprule
$A$ (m) & $T_{\rm num}$ (s) & $T_{\rm exacto}$ (s) & $T/T_0$ & $1+a^2A^2-\frac34a^4A^4$ & $v_0$ (m/s) & $x_{\max}^{\rm teor}$ (m) & $x_{\max}^{\rm num}$ (m)\\
\midrule
1 & 4.5302 & 4.5302 & 1.0099 & 1.0099 & 2 & 1.428 & 1.428 \\
4 & 5.1345 & 5.1345 & 1.1446 & 1.1408 & 4 & 2.856 & 2.856 \\
8 & 6.6283 & 6.6283 & 1.4777 & 1.333 & 6 & 4.284 & 4.284 \\
12 & 8.4782 & 8.4782 & 1.8900 & ---  & 8 & 5.711 & 5.711 \\
\bottomrule
\end{tabular}}
\caption{Izquierda: periodo en función de la amplitud (numérico frente a la fórmula elíptica \eqref{eq:Tb}). El desarrollo en serie solo vale para $aA\lesssim0.5$.
Derecha: lanzamientos desde el fondo.}\label{tab:b}
\end{table}

\begin{figure}[ht]\centering
\includegraphics[width=\textwidth]{figuras/b_periodo.pdf}
\caption{Izquierda: periodo frente a la amplitud; la simulación (círculos) coincide con la fórmula exacta \eqref{eq:Tb}, tiende a $T_0$ para amplitudes pequeñas y
a la asíntota lineal (punteada) para amplitudes grandes. Derecha: dependencia con la pendiente $a$; para $aA\ll1$ se cumple $T\propto a^{-1/2}$ y aparece un mínimo del periodo cuando $aA\sim0.5$.}
\label{fig:bper}
\end{figure}

\paragraph{Influencia de los parámetros.} La masa no interviene. El periodo escala con $g^{-1/2}$: el mismo vehículo con $A=8$~m tardaría 6.63~s en la Tierra, 16.31~s en la Luna ($g=1.62$) y 4.17~s en Júpiter ($g=24.8$) (figura~\ref{fig:benergia}, derecha).
La pendiente $a$ juega un doble papel (figura~\ref{fig:bper}, derecha): una parábola más cerrada aumenta la fuerza restauradora ($T_0\propto a^{-1/2}$), pero
también aumenta la no linealidad ($aA$), lo que alarga el periodo; el resultado es un mínimo de $T$ en $a\approx0.5/A$.

\begin{figure}[ht]\centering
\includegraphics[width=\textwidth]{figuras/b_energia_g.pdf}
\caption{Izquierda: error relativo de la energía durante 200~s. Derecha: influencia de la gravedad ($A=8$~m).}
\label{fig:benergia}
\end{figure}

% =============================================================================
\section{Apartado c): rails giratorios «El Valle Giratorio»}\label{sec:c}
% =============================================================================
\textbf{Nombre propuesto:} \emph{El Valle Giratorio}.

\paragraph{Hipótesis.} Los rails giran alrededor del eje $z$ tras recibir una velocidad angular inicial y el vehículo siempre permanece sobre ellos, luego
el ángulo azimutal del vehículo $\varphi$ coincide con el de los rails. Se supone que la inercia de los rails es despreciable
frente a la del vehículo (se relajará en la sección~\ref{sec:c:variantes}); entonces el sistema es equivalente a una partícula sobre
la superficie de revolución $z=a\rho^2$ ($\rho$ es la distancia al eje).

\paragraph{Grados de libertad y coordenadas.}
Tres coordenadas cartesianas y una ligadura holónoma, $z-a(x^2+y^2)=0$, dan $3-1=\mathbf 2$ \textbf{grados de libertad}.
Se eligen las coordenadas \emph{cilíndricas} $(q_1,q_2)=(\rho,\varphi)$: $\vec r=(\rho\cos\varphi,\ \rho\sin\varphi,\ a\rho^2)$. Dentro del plano de los rails
se puede interpretar $\rho$ como la abscisa $x$ con signo.

\paragraph{Lagrangiana.} De $v^2=\dot\rho^2+\rho^2\dot\varphi^2+\dot z^2$, $\dot z=2a\rho\dot\rho$:
\begin{equation}\label{eq:Lc}
 L_c(\rho,\dot\rho,\dot\varphi)=\tfrac12m\left[\left(1+4a^2\rho^2\right)\dot\rho^2+\rho^2\dot\varphi^2\right]-mga\rho^2 .
\end{equation}

\paragraph{Constantes del movimiento.}
\begin{itemize}[leftmargin=*]
 \item $\varphi$ es \textbf{cíclica} ($\partial L/\partial\varphi=0$, simetría de rotación en torno a $z$), luego se conserva el momento conjugado, que es la componente vertical del momento angular:
 \begin{equation}
  l\equiv p_\varphi=m\rho^2\dot\varphi=\text{cte}.
 \end{equation}
 \item $\partial L/\partial t=0$ y la ligadura es esclerónoma, luego se conserva la \textbf{energía}
 $E=\tfrac12m\left[(1+4a^2\rho^2)\dot\rho^2+\rho^2\dot\varphi^2\right]+mga\rho^2$.
\end{itemize}
Son dos constantes para dos grados de libertad: el problema es \emph{integrable por cuadraturas}.

\paragraph{Ecuaciones del movimiento.}
\begin{align}
 \text{($\varphi$):}\quad & \ddt{}\!\left(m\rho^2\dot\varphi\right)=0 \;\Longrightarrow\; \dot\varphi=\frac{l}{m\rho^2},\\[1mm]
 \text{($\rho$):}\quad & \boxed{\left(1+4a^2\rho^2\right)\ddot\rho+4a^2\rho\,\dot\rho^2-\rho\dot\varphi^2+2ga\,\rho=0}\label{eq:movc}
\end{align}
(apéndice~\ref{app:c}). Sustituyendo $\dot\varphi$, la ecuación radial queda desacoplada y es la de un problema unidimensional con masa efectiva
variable y un \emph{potencial efectivo}
\begin{equation}
 U_{\rm ef}(\rho)=\frac{l^2}{2m\rho^2}+mga\rho^2 ,\qquad E=\tfrac12m(1+4a^2\rho^2)\dot\rho^2+U_{\rm ef}(\rho).
\end{equation}

\subsection{Análisis cualitativo y casos límite}
\paragraph{Puntos de retorno.} En un retorno, $\dot\rho=0$, y $E=U_{\rm ef}$ da una ecuación cuadrática en $\rho^2$:
\begin{equation}\label{eq:retornos}
 \rho_\pm^2=\frac{E\pm\sqrt{E^2-2ga\,l^2}}{2mga}.
\end{equation}
Si $l\neq0$ la barrera centrífuga $l^2/2m\rho^2$ impide llegar al eje: $\rho\in[\rho_-,\rho_+]$, siempre acotado. Del producto de raíces:
$\rho_-\rho_+=l/(m\sqrt{2ga})$. En particular, si el vehículo parte de $\rho_0$ con $\dot\rho_0=0$ y $\dot\varphi_0=f\,\Omega_c$ con
\begin{equation}
 \Omega_c=\sqrt{2ga}=1.401\ \mathrm{rad/s}\ (13.4\ \mathrm{rpm}),
\end{equation}
el otro retorno está en $\rho=f\rho_0$ (figura~\ref{fig:cvs}, comprobado numéricamente).

\paragraph{Órbita circular.} Si $E^2=2ga\,l^2$ las dos raíces coinciden: $\rho_c^2=l/(m\sqrt{2ga})$ y
$\dot\varphi=\Omega_c=\sqrt{2ga}$, \emph{independiente del radio}. Es el equivalente a un «péndulo cónico» con
$\tan\theta=2a\rho$, y todas las órbitas circulares tienen el mismo periodo $2\pi/\Omega_c=4.49$~s (a diferencia del péndulo cónico habitual).
Es estable porque $U_{\rm ef}''(\rho_c)=8mga>0$.

\paragraph{Caso límite $l\to0$.} El vehículo pasa cada vez más cerca del eje ($\rho_-\to0$); en el límite $l=0$ el movimiento es plano y se reduce exactamente al del apartado b) (figura~\ref{fig:cl0}).

\begin{figure}[ht]\centering
\includegraphics[width=\textwidth]{figuras/c_vista_superior.pdf}
\caption{Vista superior de las trayectorias para $\rho_0=8$~m, $\dot\rho_0=0$ y distintas $\dot\varphi_0$ en unidades de $\Omega_c$.
Se ve cómo el vehículo oscila entre $f\rho_0$ y $\rho_0$ (o entre $\rho_0$ y $\rho_0/f$) mientras gira; para $f=1$ la órbita es circular.}
\label{fig:cvs}
\end{figure}

\begin{figure}[ht]\centering
\includegraphics[width=\textwidth]{figuras/c_limite_l0.pdf}
\caption{Caso límite $l\to0$: al disminuir $\dot\varphi_0$ el vehículo se acerca al eje ($\rho_{\min}=0.80$~m y $0.16$~m). Para $l=0$ (derecha) se recupera el movimiento plano del apartado b).}
\label{fig:cl0}
\end{figure}

\begin{figure}[ht]\centering
\begin{minipage}{0.48\textwidth}\centering
\includegraphics[width=\textwidth]{figuras/c_potencial_efectivo.pdf}
\caption{Potencial efectivo para las condiciones de la figura~\ref{fig:cvs}. El punto marca el estado inicial.}
\label{fig:cuef}
\end{minipage}\hfill
\begin{minipage}{0.48\textwidth}\centering
\includegraphics[width=\textwidth]{figuras/c_trayectoria_3d.png}
\caption{Trayectoria 3D sobre el paraboloide para $\dot\varphi_0=0.8\,\Omega_c$ y $\rho_0=8$ m.}
\label{fig:c3d}
\end{minipage}
\end{figure}

\subsection{Pequeñas oscilaciones radiales y precesión}
Alrededor de la órbita circular $\rho_c$ la perturbación radial oscila con
\begin{equation}
 \omega_r^2=\frac{U_{\rm ef}''(\rho_c)}{m(1+4a^2\rho_c^2)}=\frac{8ga}{1+4a^2\rho_c^2}
 \qquad\Longrightarrow\qquad \frac{\omega_r}{\Omega_c}=\frac{2}{\sqrt{1+4a^2\rho_c^2}} .
\end{equation}
Entre dos retornos consecutivos el vehículo avanza en azimut
\begin{equation}\label{eq:apsidal}
 \Delta\varphi_{\rm aps}=\pi\frac{\Omega_c}{\omega_r}=\frac\pi2\sqrt{1+4a^2\rho_c^2}.
\end{equation}
Para órbitas pequeñas $\Delta\varphi_{\rm aps}\to\pi/2$: la trayectoria es una elipse centrada en el eje (potencial armónico 2D, $z\simeq a\rho^2$), cerrada.
Para $a\rho_c\gtrsim0.5$ la órbita \emph{precesiona}: no es cerrada en general. La simulación coincide con \eqref{eq:apsidal} en todos los casos probados:

\begin{table}[ht]\centering
\begin{tabular}{rcc}
\toprule
$\rho_c$ (m) & $\Delta\varphi_{\rm aps}/\pi$ (simulación) & $\Delta\varphi_{\rm aps}/\pi$ (fórmula \eqref{eq:apsidal})\\
\midrule
1 & 0.5099 & 0.5099\\
4 & 0.6403 & 0.6403\\
8 & 0.9434 & 0.9434\\
14 & 1.4866 & 1.4866\\
\bottomrule
\end{tabular}
\caption{Ángulo apsidal para $a=0.1$~m$^{-1}$ ($\varepsilon=10^{-3}$).}
\end{table}

\subsection{Estudio sistemático de parámetros y variantes}\label{sec:c:variantes}
\paragraph{Pendiente $a$ y radio.} La figura~\ref{fig:capsidal} (izquierda) muestra $\Delta\varphi_{\rm aps}$ para tres valores de $a$: lo que importa es el producto
$a\rho_c=\tan\theta/2$, es decir, la inclinación de los rails en la órbita.
La frecuencia angular de la órbita circular $\Omega_c=\sqrt{2ga}$ aumenta con la pendiente y con $g$, y no depende de $m$.

\paragraph{Inercia de los rails.} Si los rails tienen momento de inercia $I$ respecto al eje, el Lagrangiano gana el término $\tfrac12I\dot\varphi^2$ y
$p_\varphi=(m\rho^2+I)\dot\varphi$; con $\kappa=I/m$ las ecuaciones siguen siendo \eqref{eq:movc} con $\dot\varphi=l/[m(\rho^2+\kappa)]$ (apéndice~\ref{app:c}).
Resultados (figura~\ref{fig:capsidal}, derecha, $\rho_0=8$~m, $0.8\,\Omega_c$): el rango radial pasa de $[6.40,\,8.00]$ ($\kappa=0$) a $[6.11,\,8.00]$ ($\kappa=10$~m$^2$) y
$[5.15,\,8.00]$ ($\kappa=40$~m$^2$); la barrera centrífuga $l^2/[2m(\rho^2+\kappa)]$ se hace finita en $\rho=0$, y para $\kappa=160$~m$^2$ el vehículo \emph{atraviesa el eje}
(la coordenada del plano de los rails cambia de signo). La velocidad angular de la órbita circular sigue siendo $\sqrt{2ga}$.

\paragraph{Velocidad angular impuesta.} Si un motor impusiera $\dot\varphi=\Omega$ constante la ligadura sería \emph{reónoma}, habría un solo grado de libertad
y $L=\tfrac12m[(1+4a^2x^2)\dot x^2+\Omega^2x^2]-mgax^2$, con la ecuación
$(1+4a^2x^2)\ddot x+4a^2x\dot x^2+(2ga-\Omega^2)x=0$. La energía deja de conservarse (el motor trabaja) pero se conserva
$h=\tfrac12m(1+4a^2x^2)\dot x^2-\tfrac12m\Omega^2x^2+mgax^2$. Aparece una \textbf{velocidad angular crítica} $\Omega_c=\sqrt{2ga}$:
para $\Omega<\Omega_c$ el fondo es estable y se oscila con $\omega_{\rm ef}=\sqrt{2ga-\Omega^2}$; para $\Omega>\Omega_c$ es inestable y el vehículo
es expulsado exponencialmente (tasa $\sqrt{\Omega^2-2ga}$), hacia los extremos de los rails (figura~\ref{fig:comega}). Es un caso límite de interés práctico: es la velocidad de giro máxima
segura de una atracción de este tipo ($\Omega_c\approx13.4$~rpm con $a=0.1$~m$^{-1}$).

\begin{figure}[ht]\centering
\includegraphics[width=\textwidth]{figuras/c_apsidal_kappa.pdf}
\caption{Izquierda: ángulo apsidal en función del radio de la órbita circular para tres pendientes (líneas: teoría; círculos: simulación).
Derecha: influencia del momento de inercia de los rails ($\kappa=I/m$).}
\label{fig:capsidal}
\end{figure}

\begin{figure}[ht]\centering
\includegraphics[width=\textwidth]{figuras/c_omega_impuesta.pdf}
\caption{Rails con velocidad angular impuesta. Izquierda: $\Omega<\Omega_c$, el vehículo oscila (más despacio al acercarse a $\Omega_c$). Derecha: $\Omega\geq\Omega_c$, el vehículo
sale despedido del fondo.}
\label{fig:comega}
\end{figure}

\paragraph{Conservación.} En todas las simulaciones del apartado c) la energía se conserva con un error relativo $\lesssim4\times10^{-11}$ y $l$ con error $\lesssim10^{-16}$.

% =============================================================================
\section{Apartado d): reacción normal con una lagrangiana auxiliar}\label{sec:d}
% =============================================================================
\paragraph{Lagrangiana auxiliar.} Para obtener la fuerza de ligadura se \emph{reintroduce la coordenada eliminada}: se trabaja con $(x,z)$ y se añade la ligadura
$f(x,z)=z-ax^2=0$ con un multiplicador de Lagrange $\lambda(t)$:
\begin{equation}
 L'(x,z,\dot x,\dot z,\lambda)=\tfrac12m(\dot x^2+\dot z^2)-mgz+\lambda\,(z-ax^2).
\end{equation}

\paragraph{Ecuaciones del movimiento.}
\begin{equation}\label{eq:aux}
 m\ddot x=-2a\lambda\,x,\qquad m\ddot z=-mg+\lambda,\qquad z=ax^2 .
\end{equation}
Las dos primeras salen de Euler--Lagrange en $x$ y $z$; la tercera, de $\partial L'/\partial\lambda=0$.

\paragraph{Número de incógnitas.} Hay \textbf{tres incógnitas}, $x(t)$, $z(t)$ y $\lambda(t)$, y tres ecuaciones: las dos dinámicas y la de ligadura. Hay una diferencia importante entre ellas:
$\lambda$ es una variable \emph{algebraica} (no aparece $\dot\lambda$), de modo que se obtiene derivando dos veces la ligadura, $\ddot z=2a(\dot x^2+x\ddot x)$,
y eliminando $\ddot x,\ddot z$ con las ecuaciones anteriores:
\begin{equation}\label{eq:lambda}
 \lambda=\frac{m\,(g+2a\dot x^2)}{1+4a^2x^2},
\end{equation}
mientras que, al sustituir $\lambda$ en la primera ecuación, se recupera exactamente la ecuación \eqref{eq:movb} del apartado b).

\paragraph{Significado de $\lambda$.} Los términos $\lambda\,\partial f/\partial q$ son las \emph{fuerzas generalizadas de ligadura}. Como
$\nabla f=(-2ax,\,1)$ es perpendicular a los rails,
\begin{equation}
 \vec N=\lambda\nabla f=\lambda(-2ax,\,1),\qquad N=|\vec N|=\lambda\sqrt{1+4a^2x^2}=\frac{m\,(g+2a\dot x^2)}{\sqrt{1+4a^2x^2}}.
\end{equation}
Es decir, $\lambda$ es la \textbf{componente vertical de la reacción normal} de los rails, y $N$ su módulo. Equivale a $N=m\,g\cos\theta+m\,v^2/\mathcal R$ (componente del peso más fuerza centrípeta, con $\tan\theta=2ax$ y radio de curvatura $\mathcal R=(1+4a^2x^2)^{3/2}/2a$).

\subsection{Resultados}
\begin{itemize}[leftmargin=*]
 \item En el fondo ($x=0$): $N=m(g+2av^2)$ con $v^2=2gax_0^2$, es decir
 \begin{equation}
  N_{\max}=mg\left(1+4a^2x_0^2\right).
 \end{equation}
 \item En un punto de retorno ($\dot x=0$): $N=mg/\sqrt{1+4a^2x_0^2}=mg\cos\theta_0$, como en un plano inclinado.
 \item $N>0$ en todo el recorrido: el vehículo nunca se despega de los rails.
\end{itemize}
La simulación reproduce estos valores (tabla~\ref{tab:d}). La fórmula analítica \eqref{eq:lambda} se ha comprobado de forma \emph{independiente} calculando
$\vec N=m(\ddot{\vec r}-\vec g)$ por diferenciación numérica de la trayectoria, con un error relativo máximo de $7\times10^{-8}$ (figura~\ref{fig:dnormal}, centro).

\begin{table}[ht]\centering
\begin{tabular}{rcccc}
\toprule
$x_0$ (m) & $N_{\max}/mg$ (sim.) & $1+4a^2x_0^2$ & $N(x_0)/mg$ (sim.) & $1/\sqrt{1+4a^2x_0^2}$ \\
\midrule
2 & 1.1600 & 1.1600 & 0.9285 & 0.9285\\
4 & 1.6400 & 1.6400 & 0.7809 & 0.7809\\
8 & 3.5600 & 3.5600 & 0.5300 & 0.5300\\
12 & 6.7600 & 6.7600 & 0.3846 & 0.3846\\
\bottomrule
\end{tabular}
\caption{Reacción normal en el fondo y en los extremos ($a=0.1$~m$^{-1}$). Un ocupante soltado desde $x_0=12$~m experimenta casi 7~$g$ en el fondo,
lo que fija un límite de diseño para $a\,x_0$.}\label{tab:d}
\end{table}

\begin{figure}[ht]\centering
\includegraphics[width=\textwidth]{figuras/d_normal.pdf}
\caption{Izquierda: $N/mg$ frente al tiempo para distintas posiciones iniciales (punteadas: máximo teórico). Centro: comprobación independiente de $N$ por diferenciación numérica ($x_0=8$ m).
Derecha: vector $\vec N$ a lo largo de la trayectoria; siempre perpendicular a los rails.}
\label{fig:dnormal}
\end{figure}

% =============================================================================
\section{Simulación y código}\label{sec:codigo}
% =============================================================================
Todo el código está en el repositorio \href{\repoHREF}{\texttt{\repoURL}} (Python 3; \texttt{numpy}, \texttt{scipy}, \texttt{matplotlib}, \texttt{pillow}):
\begin{center}\small
\resizebox{\textwidth}{!}{%
\begin{tabular}{ll}
\toprule
\texttt{simulacion/modelo.py} & ecuaciones, energías, $N$, periodo exacto, ángulo apsidal\dots (apéndice~\ref{app:codigo})\\
\texttt{simulacion/figuras.py} & genera todas las figuras de este informe y \texttt{resumen\_numerico.txt}\\
\texttt{simulacion/animaciones.py} & genera las animaciones GIF\\
\texttt{animaciones/b\_valle\_parabolico.gif} & movimiento sobre los rails fijos, con energías\\
\texttt{animaciones/c\_valle\_giratorio.gif} & rails giratorios en 3D y vista superior\\
\texttt{animaciones/d\_reaccion\_normal.gif} & vectores $\vec N$ y $m\vec g$ y $N(t)$\\
\bottomrule
\end{tabular}}
\end{center}
Las animaciones (también directamente en \href{\repoHREF/tree/main/animaciones}{\texttt{animaciones/}}) integran numéricamente las ecuaciones \eqref{eq:movb}, \eqref{eq:movc} y \eqref{eq:aux} (no son un dibujo del movimiento esperado).
Para reproducir todo: \texttt{pip install numpy scipy matplotlib pillow}, \texttt{python simulacion/figuras.py} y \texttt{python simulacion/animaciones.py}.

% =============================================================================
\section{Conclusiones}
% =============================================================================
\begin{enumerate}[leftmargin=*]
 \item Con rails fijos el sistema tiene un grado de libertad y es un oscilador \emph{no isócrono}: $T=T_0\,\mathcal E(-4a^2A^2)\cdot2/\pi$, que crece con la amplitud de $T_0=4.49$~s ($A\to0$) a $6.63$~s ($A=8$~m).
 \item Con los rails giratorios aparece la simetría $\varphi$ y la conservación de $l$; el vehículo queda confinado entre dos radios dados analíticamente por \eqref{eq:retornos}. Las órbitas circulares tienen
 siempre $\dot\varphi=\sqrt{2ga}$, y las órbitas casi circulares precesionan con $\Delta\varphi_{\rm aps}=\frac\pi2\sqrt{1+4a^2\rho_c^2}$.
 \item La lagrangiana auxiliar convierte la reacción normal en una incógnita más ($\lambda$); su valor máximo en el fondo es $mg(1+4a^2x_0^2)$, dato esencial para el diseño de seguridad.
 \item Las simulaciones reproducen todos los resultados analíticos con precisión de máquina, lo que valida tanto las ecuaciones como el código.
\end{enumerate}

\begin{thebibliography}{9}
\bibitem{goldstein} H.~Goldstein, C.~Poole y J.~Safko, \emph{Classical Mechanics}, 3.ª ed., Addison-Wesley (2002).
\bibitem{landau} L.~D.~Landau y E.~M.~Lifshitz, \emph{Mecánica}, Reverté.
\bibitem{taylor} J.~R.~Taylor, \emph{Classical Mechanics}, University Science Books (2005).
\bibitem{marion} S.~T.~Thornton y J.~B.~Marion, \emph{Classical Dynamics of Particles and Systems}, 5.ª ed., Brooks/Cole (2004).
\end{thebibliography}

% =============================================================================
\appendix
\section{Cálculos: de las ecuaciones de Newton a las de Lagrange}\label{app:lagrange}
% =============================================================================
\paragraph{A.1 Principio de d'Alembert.}
Sea un sistema de $N$ partículas con $m_i\ddot{\vec r}_i=\vec F_i^{(a)}+\vec f_i$. Un \emph{desplazamiento virtual} $\delta\vec r_i$ es un desplazamiento infinitesimal, arbitrario,
\emph{instantáneo} (se «congela» el tiempo y las ligaduras dependientes de $t$) y compatible con las ligaduras en ese instante. Se diferencia del real $\dd\vec r_i$ en que este último ocurre en $\dd t$ y se ve afectado por la dependencia temporal
de las ligaduras. Multiplicando cada ecuación de Newton por $\delta\vec r_i$ y sumando:
\begin{equation}
 \sum_i\left(m_i\ddot{\vec r}_i-\vec F_i^{(a)}\right)\cdot\delta\vec r_i=\sum_i\vec f_i\cdot\delta\vec r_i .
\end{equation}
Para ligaduras \emph{ideales} (sin rozamiento: $\vec f_i$ perpendicular a las superficies de ligadura y $\delta\vec r_i$ tangente a ellas), $\sum_i\vec f_i\cdot\delta\vec r_i=0$, y se obtiene el principio de d'Alembert:
\begin{equation}
 \sum_i\left(\vec F_i^{(a)}-m_i\ddot{\vec r}_i\right)\cdot\delta\vec r_i=0 .\tag{A.1}
\end{equation}
Aún no se puede igualar a cero cada término, porque los $\delta\vec r_i$ no son independientes.

\paragraph{A.2 Grados de libertad, ligaduras, coordenadas generalizadas.}
Las ligaduras holónomas son $f_\alpha(\vec r_1,\ldots,\vec r_N,t)=0$, $\alpha=1,\ldots,k$. Hay $n=3N-k$ grados de libertad. Se eligen $n$ coordenadas generalizadas independientes $q_1,\ldots,q_n$ tales que
$\vec r_i=\vec r_i(q_1,\ldots,q_n,t)$ satisface identicamente las ligaduras. Entonces
\begin{equation}
 \delta\vec r_i=\sum_{j=1}^n\pd{\vec r_i}{q_j}\delta q_j,\qquad \dot{\vec r}_i=\sum_j\pd{\vec r_i}{q_j}\dot q_j+\pd{\vec r_i}{t},
\end{equation}
con los $\delta q_j$ \emph{independientes}. De la segunda expresión se siguen las dos identidades de Lagrange:
\begin{equation}
 \pd{\dot{\vec r}_i}{\dot q_j}=\pd{\vec r_i}{q_j},\qquad \ddt{}\!\left(\pd{\vec r_i}{q_j}\right)=\pd{\dot{\vec r}_i}{q_j}.\tag{A.2}
\end{equation}

\paragraph{A.3 Ecuaciones de Lagrange.}
\emph{Trabajo virtual de las fuerzas aplicadas:} $\sum_i\vec F_i^{(a)}\cdot\delta\vec r_i=\sum_jQ_j\delta q_j$, con $Q_j=\sum_i\vec F_i^{(a)}\cdot\partial\vec r_i/\partial q_j$.

\emph{Término inercial:} usando (A.2),
\begin{align}
 \sum_im_i\ddot{\vec r}_i\cdot\pd{\vec r_i}{q_j}
 &=\ddt{}\!\left(\sum_im_i\dot{\vec r}_i\cdot\pd{\vec r_i}{q_j}\right)-\sum_im_i\dot{\vec r}_i\cdot\ddt{}\!\left(\pd{\vec r_i}{q_j}\right)\nonumber\\
 &=\ddt{}\!\left(\sum_im_i\dot{\vec r}_i\cdot\pd{\dot{\vec r}_i}{\dot q_j}\right)-\sum_im_i\dot{\vec r}_i\cdot\pd{\dot{\vec r}_i}{q_j}
 =\ddt{}\!\left(\pd{T}{\dot q_j}\right)-\pd{T}{q_j},
\end{align}
con $T=\sum_i\tfrac12m_i\dot{\vec r}_i^{\,2}$. Sustituyendo en (A.1):
$\sum_j\left[\ddt{}\!\left(\pd{T}{\dot q_j}\right)-\pd{T}{q_j}-Q_j\right]\delta q_j=0$. Como los $\delta q_j$ son independientes, cada corchete se anula:
\begin{equation}
 \ddt{}\!\left(\pd{T}{\dot q_j}\right)-\pd{T}{q_j}=Q_j .\tag{A.3}
\end{equation}
\emph{Campos conservativos:} si $\vec F_i^{(a)}=-\nabla_iV(\vec r_1,\ldots,\vec r_N,t)$, entonces $Q_j=-\sum_i\nabla_iV\cdot\partial\vec r_i/\partial q_j=-\partial V/\partial q_j$. Como $V$ no depende de las
velocidades, $\partial T/\partial\dot q_j=\partial L/\partial\dot q_j$ con $L=T-V$, y (A.3) se convierte en
\begin{equation}
 \ddt{}\!\left(\pd{L}{\dot q_j}\right)-\pd{L}{q_j}=0 .\tag{A.4}
\end{equation}

\paragraph{A.4 Relación lógica.}
La cadena es: (Newton) $+$ (ligaduras holónomas e ideales) $\Rightarrow$ (A.1) $\Rightarrow$ (A.3) $\Rightarrow$ (A.4) si las fuerzas son conservativas.
Cada paso es reversible: de (A.4) se llega a (A.3) y, como $\partial\vec r_i/\partial q_j$ tiene rango máximo, a (A.1), que es la proyección de las ecuaciones de Newton sobre las direcciones compatibles con las ligaduras.
Para una partícula libre en cartesianas, (A.4) \emph{es} la segunda ley de Newton ($m\ddot x_j=-\partial V/\partial x_j$). Newton es el postulado de partida y Lagrange una reformulación equivalente
(con la hipótesis adicional de ligaduras ideales, que no está en las leyes de Newton).

% =============================================================================
\section{Cálculos del apartado b)}\label{app:b}
% =============================================================================
\paragraph{B.1 Lagrangiana y ecuación del movimiento.}
Con $z=ax^2$, $\dot z=2ax\dot x$: $T=\tfrac12m\dot x^2(1+4a^2x^2)$, $V=mgax^2$. Entonces
\begin{align}
 \pd{L}{\dot x}&=m(1+4a^2x^2)\dot x,&
 \ddt{}\!\pd{L}{\dot x}&=m(1+4a^2x^2)\ddot x+8ma^2x\dot x^2,\\
 \pd{L}{x}&=4ma^2x\dot x^2-2mgax .
\end{align}
La ecuación de Euler--Lagrange da $m(1+4a^2x^2)\ddot x+8ma^2x\dot x^2-4ma^2x\dot x^2+2mgax=0$, es decir,
\begin{equation}
 (1+4a^2x^2)\ddot x+4a^2x\dot x^2+2gax=0.
\end{equation}

\paragraph{B.2 Energía y función de Jacobi.}
$h=\dot x\,\partial L/\partial\dot x-L=m(1+4a^2x^2)\dot x^2-\tfrac12m(1+4a^2x^2)\dot x^2+mgax^2=\tfrac12m(1+4a^2x^2)\dot x^2+mgax^2=E$.
$\dd h/\dd t=-\partial L/\partial t=0$.

\paragraph{B.3 Periodo exacto.}
De $E=mgaA^2$ ($A$: punto de retorno, $\dot x=0$):
\begin{equation}
 \dot x=\pm\sqrt{\frac{2ga(A^2-x^2)}{1+4a^2x^2}} .
\end{equation}
El tiempo en ir de $x=0$ a $x=A$ es un cuarto del periodo:
\begin{equation}
 T=4\int_0^A\sqrt{\frac{1+4a^2x^2}{2ga\,(A^2-x^2)}}\,\dd x .
\end{equation}
Con $x=A\sin\theta$ ($\dd x=A\cos\theta\,\dd\theta$, $\sqrt{A^2-x^2}=A\cos\theta$):
\begin{equation}
 T=\frac{4}{\sqrt{2ga}}\int_0^{\pi/2}\sqrt{1+4a^2A^2\sin^2\theta}\,\dd\theta=\frac{4}{\sqrt{2ga}}\,\mathcal E(-4a^2A^2).
\end{equation}
Desarrollando $\mathcal E(m)=\frac\pi2\left[1-\frac m4-\frac{3m^2}{64}-\dots\right]$ con $m=-4a^2A^2$:
$T=T_0\left[1+a^2A^2-\tfrac34a^4A^4+\dots\right]$, con $T_0=2\pi/\sqrt{2ga}$. Para $aA\gg1$ el integrando es $\simeq2aA\sin\theta$, $\int_0^{\pi/2}\sin\theta\,\dd\theta=1$,
y $T\simeq8aA/\sqrt{2ga}=4\sqrt{2a/g}\,A$.

\paragraph{B.4 Pequeñas oscilaciones.} Para $|x|\ll1/(2a)$ y velocidades pequeñas se desprecian los términos $4a^2x^2\ddot x$ y $4a^2x\dot x^2$ (de orden superior): $\ddot x+2ga\,x=0$, $\omega_0^2=2ga$.
Si se lanza desde el fondo con $v_0$: $\tfrac12mv_0^2=mgax_{\max}^2$, $x_{\max}=v_0/\sqrt{2ga}$.

% =============================================================================
\section{Cálculos del apartado c)}\label{app:c}
% =============================================================================
\paragraph{C.1 Lagrangiana.}
$\vec r=(\rho\cos\varphi,\rho\sin\varphi,a\rho^2)$, $\dot{\vec r}=(\dot\rho\cos\varphi-\rho\dot\varphi\sin\varphi,\ \dot\rho\sin\varphi+\rho\dot\varphi\cos\varphi,\ 2a\rho\dot\rho)$,
$v^2=\dot\rho^2+\rho^2\dot\varphi^2+4a^2\rho^2\dot\rho^2$. $L=\tfrac12mv^2-mga\rho^2$ da \eqref{eq:Lc}.

\paragraph{C.2 Ecuaciones.}
\begin{align}
 \varphi:&\ \ \ddt{}\!\left(m\rho^2\dot\varphi\right)=0,\\
 \rho:&\ \ \ddt{}\!\left[m(1+4a^2\rho^2)\dot\rho\right]-\left[4ma^2\rho\dot\rho^2+m\rho\dot\varphi^2-2mga\rho\right]=0\nonumber\\
 &\Longrightarrow\ (1+4a^2\rho^2)\ddot\rho+8a^2\rho\dot\rho^2-4a^2\rho\dot\rho^2-\rho\dot\varphi^2+2ga\rho=0,
\end{align}
que es \eqref{eq:movc}. Con $\dot\varphi=l/m\rho^2$ el sistema se reduce a una ecuación para $\rho(t)$ y a una cuadratura para $\varphi(t)$.

\paragraph{C.3 Energía y potencial efectivo.}
$h=\dot\rho\,p_\rho+\dot\varphi\,p_\varphi-L=E$, con $E=\tfrac12m(1+4a^2\rho^2)\dot\rho^2+\tfrac12m\rho^2\dot\varphi^2+mga\rho^2$. Usando $\dot\varphi=l/m\rho^2$,
$E=\tfrac12m(1+4a^2\rho^2)\dot\rho^2+U_{\rm ef}$, $U_{\rm ef}=\frac{l^2}{2m\rho^2}+mga\rho^2$.

\paragraph{C.4 Puntos de retorno.}
En $\dot\rho=0$: $E=\frac{l^2}{2m\rho^2}+mga\rho^2$. Con $u=\rho^2$: $mga\,u^2-E\,u+\frac{l^2}{2m}=0$, cuyas raíces son
$u_\pm=\frac{E\pm\sqrt{E^2-2ga\,l^2}}{2mga}$ (hay órbita acotada si $E^2\geq2ga\,l^2$). Producto de raíces: $u_+u_-=\frac{l^2}{2m^2ga}$ $\Rightarrow$ $\rho_+\rho_-=\frac{l}{m\sqrt{2ga}}$.
Si parte de un retorno $\rho_0$ con $\dot\varphi_0$: $l/m=\rho_0^2\dot\varphi_0$ y $\rho_{\rm otro}=\rho_0\dot\varphi_0/\sqrt{2ga}=f\rho_0$.

\paragraph{C.5 Órbita circular y pequeñas oscilaciones.}
$U_{\rm ef}'=-\frac{l^2}{m\rho^3}+2mga\rho=0\Rightarrow\rho_c^4=\frac{l^2}{2m^2ga}$. Entonces $\dot\varphi_c=\frac{l}{m\rho_c^2}=\sqrt{2ga}$ (independiente de $\rho_c$).
$U_{\rm ef}''=\frac{3l^2}{m\rho^4}+2mga=6mga+2mga=8mga$ en $\rho_c$. Linealizando $\rho=\rho_c+\xi$ (en la ecuación radial el término $4a^2\rho\dot\rho^2$ es de segundo orden):
$m(1+4a^2\rho_c^2)\ddot\xi=-U_{\rm ef}''\xi\Rightarrow\omega_r^2=\frac{8ga}{1+4a^2\rho_c^2}$. Entre dos retornos transcurre $\pi/\omega_r$, y en ese tiempo el vehículo gira $\Omega_c\pi/\omega_r=\frac\pi2\sqrt{1+4a^2\rho_c^2}$.

\paragraph{C.6 Inercia de los rails ($\kappa=I/m$).}
Añadiendo $\tfrac12I\dot\varphi^2$ a $L_c$: $p_\varphi=m(\rho^2+\kappa)\dot\varphi$ constante; la ecuación radial es la misma \eqref{eq:movc} con $\dot\varphi=L/(\rho^2+\kappa)$, $L=l/m$;
$U_{\rm ef}/m=\frac{L^2}{2(\rho^2+\kappa)}+ga\rho^2$. Circular: $(\rho^2+\kappa)^2=L^2/2ga\Rightarrow\dot\varphi_c=\sqrt{2ga}$ de nuevo, y
$\omega_r^2=\dfrac{8ga\rho_c^2}{(\rho_c^2+\kappa)(1+4a^2\rho_c^2)}$.

\paragraph{C.7 Velocidad angular impuesta.}
Con $\varphi=\Omega t$, $\vec r=(x\cos\Omega t,x\sin\Omega t,ax^2)$, $L=\tfrac12m[(1+4a^2x^2)\dot x^2+\Omega^2x^2]-mgax^2$.
Euler--Lagrange: $(1+4a^2x^2)\ddot x+4a^2x\dot x^2+(2ga-\Omega^2)x=0$. $L$ no depende de $t$ pero $L\neq T-V$
(aparece $\Omega^2x^2/2$ en $T_2$ sin velocidad), luego $h=\dot x\,p_x-L=\tfrac12m(1+4a^2x^2)\dot x^2-\tfrac12m\Omega^2x^2+mgax^2$ se conserva pero $\neq E$.
Linealizando: $\ddot x=-(2ga-\Omega^2)x$.

% =============================================================================
\section{Cálculos del apartado d)}\label{app:d}
% =============================================================================
\paragraph{D.1 Multiplicadores de Lagrange.}
Para ligaduras holónomas $f_\alpha(q,t)=0$ en $n$ coordenadas «sobrantes», $L'=L+\sum_\alpha\lambda_\alpha f_\alpha$ y
$\ddt{}\pd{L}{\dot q_j}-\pd{L}{q_j}=\sum_\alpha\lambda_\alpha\pd{f_\alpha}{q_j}$ ($=Q^{\rm lig}_j$). Aquí $q=(x,z)$, $f=z-ax^2$:
\begin{align}
 x:&\ \ m\ddot x=-2a\lambda x,\\
 z:&\ \ m\ddot z=-mg+\lambda,\\
 \lambda:&\ \ z-ax^2=0 .
\end{align}

\paragraph{D.2 Cálculo de $\lambda$.}
De la ligadura, $\dot z=2ax\dot x$, $\ddot z=2a(\dot x^2+x\ddot x)$. Sustituyendo $\ddot x=-2a\lambda x/m$:
$\ddot z=2a\dot x^2-4a^2\lambda x^2/m$. Igualando a $-g+\lambda/m$:
\begin{equation}
 \frac{\lambda}{m}\left(1+4a^2x^2\right)=g+2a\dot x^2\ \Rightarrow\ \lambda=\frac{m(g+2a\dot x^2)}{1+4a^2x^2}.
\end{equation}
Sustituyendo en la ecuación de $x$: $\ddot x=-\frac{2ax(g+2a\dot x^2)}{1+4a^2x^2}$, es decir, $(1+4a^2x^2)\ddot x+4a^2x\dot x^2+2gax=0$ (ec.~\eqref{eq:movb}).

\paragraph{D.3 Reacción normal.}
La fuerza de ligadura es $\vec N=(Q^{\rm lig}_x,Q^{\rm lig}_z)=\lambda(-2ax,\,1)=\lambda\nabla f$, perpendicular a los rails. $N=\lambda\sqrt{1+4a^2x^2}=\dfrac{m(g+2a\dot x^2)}{\sqrt{1+4a^2x^2}}$.
En el fondo: $x=0$, $\dot x^2=2gax_0^2$: $N=m(g+4ga^2x_0^2)=mg(1+4a^2x_0^2)$. En un retorno: $N=mg/\sqrt{1+4a^2x_0^2}=mg\cos\theta_0$.
Comprobación con la geometría: $\tan\theta=2ax$, $\cos\theta=1/\sqrt{1+4a^2x^2}$, $v^2=(1+4a^2x^2)\dot x^2$, curvatura $2a/(1+4a^2x^2)^{3/2}$:
$N=mg\cos\theta+mv^2\cdot\frac{2a}{(1+4a^2x^2)^{3/2}}=\frac{m(g+2a\dot x^2)}{\sqrt{1+4a^2x^2}}$, que coincide.

% =============================================================================
\section{Código del modelo}\label{app:codigo}
% =============================================================================
Listado completo de \texttt{simulacion/modelo.py} (el resto del código, que solo dibuja, está en el repositorio).
\lstinputlisting{../simulacion/modelo.py}

\end{document}
